%% =========================================================================
\section{Conclusion}
%% =========================================================================

This report looked at synchronous generator behaviour under three increasingly
complex scenarios using MATLAB R2025b / Simscape Electrical. The results were
checked against standard power systems equations.

\textbf{Exercise~1} covered the two main control loops of a synchronous
machine: the swing-equation/governor loop linking active power to frequency,
and the Q-equation/AVR loop linking reactive power to terminal voltage. The
\MVA{555} machine recovered from a \pu{0.82} voltage dip within 3--4 seconds
thanks to AVR action, and the frequency settled at values matching the droop
prediction from \eqnref{eq:droop_calc_3s}.

\textbf{Exercise~2} showed that the governor frequency setpoint controls load
sharing between parallel generators. When the oncoming machine has a lower
setpoint (Task~3), it draws power from the system, which is a dangerous
motoring condition. A higher setpoint (Task~4) makes the new machine take over
most of the load, which is the intended outcome. In both cases, the system
frequency settled at a droop-weighted average of the two setpoints, consistent
with \eqnref{eq:combined_droop}.

\textbf{Exercise~3} showed how an infinite bus differs from a finite one.
Connected to an infinite bus, Generator~B has no effect on system frequency.
Its power output is set entirely by its governor setpoint relative to the fixed
grid frequency --- Task~5 resulted in motoring and Task~6 gave a clean power
injection. In practice, this is what happens when a small generator connects to
a national grid. The grid frequency is set by the combined inertia of millions
of megawatts, and one machine cannot change it.
